% =====================================================================
%  PREAMBOLO - Dispensa Programmazione 2 (UNICT)
%  Compilare main.tex con pdflatex (2 volte per l'indice).
% =====================================================================

% ---------- Lingua, font, codifica ----------
\usepackage[utf8]{inputenc}
\usepackage[T1]{fontenc}
\usepackage[italian]{babel}
\usepackage{lmodern}
\usepackage{microtype}
\DisableLigatures{encoding = T1, family = tt*}   % evita che >> diventi » nel codice
\usepackage{textcomp}

% ---------- Pagina ----------
\usepackage[a4paper,margin=2.3cm,headheight=14pt]{geometry}
\usepackage{fancyhdr}
\pagestyle{fancy}
\fancyhf{}
\renewcommand{\chaptermark}[1]{\markboth{#1}{}}
\fancyhead[L]{\small\nouppercase{\leftmark}}
\fancyfoot[L]{\small Programmazione 2 -- UNICT}
\fancyfoot[R]{\thepage}
\setlength{\parindent}{0pt}
\setlength{\parskip}{0.55em}

% ---------- Matematica ----------
\usepackage{amsmath}
\usepackage{amssymb}

% ---------- Colori: versione LIGHT (stampa) o DARK (schermo) ----------
% dispensa-dark.tex definisce \modalitadark prima di caricare il preambolo.
% Trucco per il dark: si ridefiniscono "white" (= sfondo) e "black" (= testo).
% Così tutte le miscele tipo blue!8!sfondo (8% blu + 92% "white") diventano
% automaticamente tinte scure, e red!70!black diventa un rosso chiaro.
\usepackage[dvipsnames,table]{xcolor}
\ifdefined\modalitadark
  \definecolor{sfondo}{HTML}{1E1F22}   % usa SEMPRE "sfondo" nelle tinte: x!8!sfondo
  \definecolor{black}{HTML}{E3E3E3}    % il testo (e le miscele x!y!black) diventano chiari
  \definecolor{codebg}{HTML}{27292E}
  \definecolor{codekw}{HTML}{6CB6FF}
  \definecolor{codecom}{HTML}{7EC699}
  \definecolor{codestr}{HTML}{F0A35E}
  \definecolor{codenum}{HTML}{7A7A7A}
  \definecolor{stackcol}{HTML}{23384A}
  \definecolor{heapcol}{HTML}{4A3B22}
  \definecolor{privcol}{HTML}{4D2A2A}
  \definecolor{pubcol}{HTML}{223F2E}
  % il colore del testo diventa anche quello "di default" (header, footer, numeri)
  \makeatletter
  \AtBeginDocument{\pagecolor{sfondo}\color{black}\let\default@color\current@color}
  \makeatother
\else
  \definecolor{sfondo}{RGB}{255,255,255}
  \definecolor{codebg}{RGB}{248,248,250}
  \definecolor{codekw}{RGB}{0,72,170}
  \definecolor{codecom}{RGB}{0,128,64}
  \definecolor{codestr}{RGB}{190,80,0}
  \definecolor{codenum}{RGB}{140,140,140}
  \definecolor{stackcol}{RGB}{214,234,248}
  \definecolor{heapcol}{RGB}{253,235,208}
  \definecolor{privcol}{RGB}{250,219,216}
  \definecolor{pubcol}{RGB}{213,245,227}
\fi

% ---------- Tabelle ed elenchi ----------
\usepackage{booktabs}
\usepackage{array}
\usepackage{tabularx}
\usepackage{enumitem}
\usepackage{needspace}   % \needspace{8cm} prima di "titoletto + figura": li tiene sulla stessa pagina
\setlist{itemsep=0.25em,topsep=0.3em}

% ---------- Codice C++ (listings) ----------
\usepackage{listings}
\lstdefinestyle{cpp}{
  language=C++,
  basicstyle=\ttfamily\small,
  keywordstyle=\color{codekw}\bfseries,
  commentstyle=\color{codecom}\itshape,
  stringstyle=\color{codestr},
  numberstyle=\tiny\color{codenum},
  numbers=left,
  numbersep=8pt,
  backgroundcolor=\color{codebg},
  frame=single,
  rulecolor=\color{black!25!sfondo},
  framesep=4pt,
  xleftmargin=1.6em,
  breaklines=true,
  postbreak=\mbox{\textcolor{red}{$\hookrightarrow$}\space},
  showstringspaces=false,
  tabsize=4,
  columns=fullflexible,
  keepspaces=true,
  morekeywords={constexpr,nullptr,override,string,namespace,using,cout,cin,cerr,endl,getline,stringstream},
  literate=
    {à}{{\`a}}1 {è}{{\`e}}1 {é}{{\'e}}1 {ì}{{\`i}}1
    {ò}{{\`o}}1 {ù}{{\`u}}1 {È}{{\`E}}1 {°}{{\textdegree}}1
    {À}{{\`A}}1 {É}{{\'E}}1 {Ì}{{\`I}}1 {Ò}{{\`O}}1 {Ù}{{\`U}}1
}
\lstset{style=cpp}
\renewcommand{\lstlistingname}{Codice}
% Codice inline: \cc|cin >> x|
\newcommand{\cc}{\lstinline[basicstyle=\ttfamily\normalsize,keywordstyle=\color{codekw},postbreak={}]}
\newcommand{\nl}{\texttt{\textbackslash n}}
% Output REALE di un programma (file generati da verifica_codice.sh)
\newcommand{\uscita}[1]{\lstinputlisting[language={},numbers=none,
  backgroundcolor=\color{black!6!sfondo},basicstyle=\ttfamily\small,
  title={\footnotesize\itshape Output reale del programma}]{#1}}
% Input dato al programma (file codice/mNN/input/X.txt)
\newcommand{\ingresso}[1]{\lstinputlisting[language={},numbers=none,
  backgroundcolor=\color{yellow!8!sfondo},basicstyle=\ttfamily\small,
  title={\footnotesize\itshape Input digitato (non compare nell'output)}]{#1}}
% Codice da file con didascalia = percorso
\newcommand{\codice}[1]{\lstinputlisting[caption={\texttt{\detokenize{#1}}}]{#1}}
\newcommand{\q}[1]{\textquotedbl#1\textquotedbl}   % virgolette dritte sicure

% ---------- TikZ ----------
\usepackage{graphicx}
\usepackage{tikz}
\usetikzlibrary{positioning,arrows.meta,shapes.multipart,shapes.misc,calc,fit,backgrounds,decorations.pathreplacing,babel}
\tikzset{
  freccia/.style={-{Stealth[length=2.6mm]},thick},
  frecciarossa/.style={-{Stealth[length=2.6mm]},thick,red!75!black},
  cella/.style={draw,minimum width=1.6cm,minimum height=0.75cm,font=\ttfamily\small},
  cellina/.style={draw,minimum width=0.75cm,minimum height=0.75cm,font=\ttfamily\small},
  etichetta/.style={font=\footnotesize\itshape,text=black!70!sfondo},
  zona/.style={draw,rounded corners,inner sep=8pt},
  classe/.style={draw,rectangle split,rectangle split parts=3,
                 rectangle split part fill={blue!15!sfondo,privcol,pubcol},
                 text width=5.2cm,font=\ttfamily\small,align=left},
  processo/.style={draw,rounded corners,fill=blue!8!sfondo,minimum height=0.9cm,
                   text width=2.4cm,align=center,font=\small},
  file/.style={draw,fill=yellow!20!sfondo,minimum height=0.9cm,
               text width=1.9cm,align=center,font=\small\ttfamily},
}

% ---------- Box didattici (tcolorbox) ----------
\usepackage[most]{tcolorbox}
\tcbset{enhanced,breakable,boxrule=0.8pt,arc=2mm,coltitle=sfondo,coltext=black,
        fonttitle=\bfseries,left=6pt,right=6pt,top=4pt,bottom=4pt}
\newtcolorbox{analogia}[1][]{colback=blue!4!sfondo,colframe=blue!55!black,
  title={Analogia#1}}
\newtcolorbox{definizione}[1][]{colback=black!3!sfondo,colframe=black!65!sfondo,
  title={Definizione#1}}
\newtcolorbox{trappola}[1][]{colback=red!5!sfondo,colframe=red!70!black,
  title={Trappola d'esame#1}}
\newtcolorbox{consiglio}[1][]{colback=green!5!sfondo,colframe=green!45!black,
  title={Tip#1}}
\newtcolorbox{orale}[1][]{colback=violet!4!sfondo,colframe=violet!60!black,
  title={#1}}
\newtcolorbox{soluzione}[1][]{colback=yellow!6!sfondo,colframe=orange!70!black,
  title={Soluzione#1}}
\newtcolorbox{traccia}[1][]{colback=sfondo,colframe=black!80!sfondo,
  title={Traccia#1}}
\newtcolorbox{passopasso}[1][]{colback=cyan!4!sfondo,colframe=cyan!50!black,
  title={Passo per passo#1}}
\newtcolorbox{inparole}[1][]{colback=orange!5!sfondo,colframe=orange!60!black,
  title={In parole povere#1}}
% Laboratorio a step: \begin{step}{titolo dello step}
\newtcolorbox[auto counter,number within=chapter,
  number freestyle={\noexpand\arabic{\tcbcounter}}]{step}[1]{%
  colback=black!2!sfondo,colframe=teal!60!black,
  title={Step \thetcbcounter\ -- #1}}
% Esercizi al PC: \begin{esercizio}{titolo}{facile|medio|difficile}
\newtcolorbox[auto counter,number within=chapter]{esercizio}[2]{%
  colback=sfondo,colframe=blue!45!black,
  title={Esercizio \thetcbcounter\ -- #1\hfill\normalfont\itshape\small difficoltà: #2}}


% ---------- Quiz a crocette ----------
\newcounter{domanda}[chapter]
\newenvironment{quiz}[1]{%
  \refstepcounter{domanda}%
  \par\medskip\noindent\textbf{Domanda \thedomanda.} #1
  \begin{enumerate}[label=\textbf{\Alph*)},leftmargin=2.2em,itemsep=0.1em]}%
  {\end{enumerate}}

% ---------- Collegamenti (sempre per ultimo) ----------
\usepackage[hidelinks]{hyperref}
\hypersetup{pdftitle={Programmazione 2 - Dispensa d'esame},pdfauthor={Giulio Casella}}
